\chapter{问题价值与作品边界}
\label{chap:value-boundary}

\section{从进程抽象到智能体工作流}

经典进程抽象回答“谁在执行、使用了多少 CPU 和内存、何时阻塞”，但智能体工作流
还要持续回答四个问题：当前决策属于哪个角色和 workflow；它依据了哪一轮上下文；
工具副作用是否在授权范围内；跨 Agent 结果是可信控制还是不可信数据。若这些信息只
存在于 prompt、JSON 或应用日志中，内核无法在真正发生文件、IPC 和资源效果之前
检查它们。

我们不试图让操作系统代替模型思考。我们把问题重新表述为：操作系统能否提供一组
足够小、能够强制、可观测又能被不同策略复用的机制，让用户态 Agent Loop 在失败、
并发和攻击输入下仍保持清晰边界。这个表述避免了两个极端。一端是把现有脚本换成
“Agent”名字却仍然硬编码所有结果；另一端是把模型 SDK、网络协议和自然语言解析
塞进教学内核，扩大可信计算基并破坏可验证性。

\begin{designbox}[我们的价值判断]
我们以可重验的系统属性评价 AgentOS：用户态智能体完成真实任务时，身份、权限、
资源、因果和证据都能由内核独立检查。
\end{designbox}

\section{需要进入内核的机制}

我们按“效果发生点”选择内核增量。进程成为 Agent 时，内核发布 identity、role、
control id 和 capability；进入 workflow 时，内核绑定不可混淆的
\code{id+generation} 与 VFS scope；工具调用时，内核解析版本化参数并检查 schema、
capability 和数据来源；消息与文件结果进入下一轮前，内核保守传播 provenance；
等待模型或事件时，内核使线程真实休眠；调度与资源效果则按 workflow 聚合，而不是
让一个 workflow 通过增加线程获得额外公平份额。

表~\ref{tab:kernel-user-boundary} 汇总可信边界。我们刻意让用户态保持丰富、内核接口
保持紧凑。这样既能使用真实模型，也能用 deterministic replay 或固定策略验证同一
底层路径。

\begin{table}[htbp]
\centering
\caption{机制与策略的责任分界}
\label{tab:kernel-user-boundary}
\begin{tabularx}{\textwidth}{L{2.5cm}YY}
\toprule
\textbf{层次} & \textbf{负责内容} & \textbf{明确不负责} \\
\midrule
AgentOS 内核 & 身份、role/capability、scope、Context、provenance、typed ABI、
route/wait、资源、调度、审计和 fence & 自然语言规划、prompt 内容判断、工具选择、
模型推理、JSON/HTTP/TLS/MCP wire \\
Guest 用户态 & 目标分解、角色策略、历史、工具呈现、TASK 状态机、artifact 组织、
错误重规划 & 绕过内核 capability、伪造 Context 因果、替换内核资源账本 \\
Host & CLI、observer、本地 socket、QEMU 串口、API key、TLS、provider 协议转换 &
读取 Guest 业务文件、选择业务工具、代写工具结果、创建内核证据 \\
\bottomrule
\end{tabularx}
\end{table}

\section{增量边界：保留 uCore 的教学可读性}

作品建立在 LearningOS/uCore 教学内核之上\cite{ucore2025}。我们保留原有进程、页表、
VFS、trap、调度和系统调用主干，只在明确边界增加 Agent 字段、workflow lifecycle、
Context 映射、工具协议、元数据、调度账户和观测设施。普通进程仍可运行；Agent 身份
只能由受控 create/exec 路径发布。所有新增 UAPI 使用定长结构、显式 version 和 size，
并以 kernel/user 共用头文件及 layout checker 固定 ABI。

我们没有修改文件系统磁盘格式。Agent metadata 是当前启动周期内显式登记的内存
目录，并绑定真实 inode incarnation；live query 只索引这些对象。我们也没有扩大
uCore 的 \code{MAXFILE} 来容纳产品程序，而是在构建期约束 Nexus Guest raw binary
不超过 274432 字节。这个取舍使新增能力仍受原教学内核资源边界约束。

\section{三类工具执行路径为何必须分开}

作品同时保留 compact batch、typed V2、ENFORCE V3 与 Task SQ/CQ，分别服务同步操作、
探索式调用、冻结合同和有界提交/完成。

\begin{table}[htbp]
\centering
\caption{工具执行路径的用途与保证}
\label{tab:tool-paths}
\begin{tabularx}{\textwidth}{L{2.3cm}L{3.0cm}YY}
\toprule
\textbf{路径} & \textbf{输入形态} & \textbf{适用场景} & \textbf{边界} \\
\midrule
Compact batch & 最多 64 个 \code{agent_op} & 固定顺序、低入口调用次数的同步操作 &
不是 typed KV、并行或 non-blocking batch \\
Typed V2 & \code{uint64}/string KV & 模型依据上一结果探索式选择下一工具 &
逐次检查，但没有冻结 DAG 的 predecessor/attempt 包络 \\
ENFORCE V3 & V2 prefix 加 contract/node/attempt/manifest & 预先声明的高保证 DAG &
动态 Agent Loop 不能假装调用图已经冻结 \\
Task SQ/CQ & 16-slot single-issuer ring & 有界提交/完成、backpressure、resync 研究 &
当前同步 null/NONE provider，无业务 payload backend \\
\bottomrule
\end{tabularx}
\end{table}

Nexus 选择 V2，使 Coordinator 能根据刚返回的 System 或 Research 工件改变下一步；
高保证固定流程使用 V3。Task Channel 当前范围止于 null/NONE provider 和有界 SQ/CQ
语义，业务 payload backend 留待后续实现。

\section{从参考思想到本项目增量}

我们参考 Linux cpuacct/rstat 的批量资源记账思想\cite{linuxcpuacct,linuxrstat}、BPF
ring buffer 的有序提交思想\cite{bpfring}、Linux EEVDF 的 lag 与 virtual deadline
概念\cite{eevdf}、io\_uring 的 SQ/CQ 通信模型\cite{iouring}、Haiku BFS 的属性和
选择性索引\cite{haikubfs}。在 Agent 方向，我们参考 tool-use 循环中“模型选择、应用
执行、结果回灌”的职责划分\cite{anthropictools,claudeloop}，并研究 MCP/A2A 的对象
形状\cite{mcp2026,a2av1}。

这些工作提供概念来源；兼容范围以本项目已经实现的协议和测试为准。我们的增量是把
相关思想约束到 uCore 的 workflow identity、Context、capability、provenance、VFS 和
evidence 语义中。Credit Domain 为同一 workflow 的 exec/storage 账户建立 U/P/F 状态和
fence 精确结算；Evidence Ring 密封 Agent Context、gap、credit 和 metadata generation；
Task Channel 当前实现 AgentOS 自有的有界 SQ/CQ ABI。

\begin{boundarybox}[原创与互操作边界]
仓库没有 vendoring Linux、Haiku、AIOS、相关论文原型、MCP/A2A SDK 或 Wasm runtime
源码。MCP/A2A 当前是 deterministic in-memory 对象映射研究，不是 server、streaming
transport、内核 adapter 或跨实现互操作证明。
\end{boundarybox}

\section{产品价值：把“Agent 想做什么”和“系统允许什么”放在同一现场}

确定性 \code{contest-demo} 解决性能复验问题，AgentOS Nexus 解决产品交互和机制
可见性问题。前者保持同一业务合同和 AB/BA 次序，不让模型随机性污染对比；后者让
评委临时提出目标，观察模型如何委派，又能在第二窗口按 task/correlation 对齐内核
route、WAITING、wake、capability、Context 与拒绝原因。

这种双产品面比单一 Dashboard 更有解释力。左侧结果说明策略确实消费了工具和工件；
右侧观测说明这些动作没有绕过操作系统。两者共同展示“用户态可自由演进，内核边界
保持可重验”的系统价值。

\section{本章验证落点}

本章边界由以下仓库事实交叉约束：公开 UAPI layout checker 固定结构
尺寸和偏移；模块 checker 限制职责归属；Guest 场景验证系统调用行为；Host tests 验证
frame、sequence、hash 和 fail-closed；one-shot campaign 冻结性能原始数据。第
\ref{chap:validation}章给出完整分层验证矩阵，第\ref{chap:limitations}章集中列出我们
没有交付的能力。
